\section{Organizzazione del lavoro}
\subsection{Suddivisione dei compiti}
\subsubsection{Database}
\subsubsection{Mockup delle pagine}
\subsubsection{Implementazione delle pagine}
\subsection{Periodo di sviluppo}
\subsection{Strumenti utilizzati}
Per il lavoro collaborativo abbiamo utilizzato \textbf{Git} come sistema di controllo versione e \textbf{GitHub} come piattaforma di hosting del repository remoto.\\
Abbiamo ritenuto opportuno utilizzare l'ambiente di rilascio per lo sviluppo de progetto per assicurarci che il sito web funzionasse correttamente nell'ambiente in cui sarebbe stato valutato.\\
Sono stati utili a questo scopo i seguenti strumenti:
\begin{itemize}
    \item \textbf{PUTTY}: client SSH per connettersi al server remoto;
    \item \textbf{Estensione SFTP per visual studio code}: per trasferire i file modificati sul server remoto automaticamente ad ogni salvataggio. L'estensione utilizzata è identificata tramite l'ID \texttt{natizyskunk.sftp} v. 1.16.3.
\end{itemize}